% Standalone research entry 207; batch 203--207.
% Original section material: Pasted text(3).txt, source edition 22 September 2026.
% Research additions: 27 September 2026.
% Compile twice with: lualatex 207.tex
% !TeX program = lualatex
% Compile twice: lualatex open_problems_500.tex
% All references and prose are embedded; no BibTeX database or external inputs.
\documentclass[11pt,a4paper]{article}
\usepackage[margin=24mm,headheight=14pt]{geometry}
\usepackage{fontspec}
\setmainfont{Latin Modern Roman}
\setsansfont{Latin Modern Sans}
\setmonofont{Latin Modern Mono}
\usepackage{amsmath,amssymb,mathtools}
\usepackage{unicode-math}
\setmathfont{Latin Modern Math}
\usepackage{microtype}
\usepackage{enumitem,xurl,needspace}
\usepackage[dvipsnames]{xcolor}
\usepackage{hyperref}
\hypersetup{unicode=true,colorlinks=true,linkcolor=MidnightBlue,urlcolor=MidnightBlue,
 pdftitle={500 Mathematical Problem Statements},pdfauthor={Source-annotated compilation},bookmarksnumbered=true}
\usepackage{bookmark}
\usepackage{tocloft}
\setlength{\cftsecnumwidth}{3em}
\cftsetpnumwidth{2.5em}
\cftsetrmarg{4em}
\usepackage{titlesec}
\titleformat{\section}{\Large\bfseries\raggedright}{\thesection}{0.7em}{}
\usepackage{fancyhdr}
\pagestyle{fancy}\fancyhf{}
\fancyhead[L]{\small\sffamily Mathematical problem statements}
\fancyhead[R]{\small Source edition 22}
\fancyfoot[C]{\thepage}
\setlength{\parindent}{0pt}
\setlength{\parskip}{5pt plus 1pt}
\setlength{\emergencystretch}{3em}
\setcounter{tocdepth}{1}
\setcounter{secnumdepth}{1}
\setlist[itemize]{leftmargin=3em,itemsep=5pt,topsep=4pt,parsep=0pt}
\allowdisplaybreaks
\newcommand{\N}{\mathbb{N}}
\newcommand{\Z}{\mathbb{Z}}
\newcommand{\Q}{\mathbb{Q}}
\newcommand{\R}{\mathbb{R}}
\newcommand{\C}{\mathbb{C}}
\newcommand{\F}{\mathbb{F}}
\newcommand{\A}{\mathbb{A}}
\newcommand{\PP}{\mathbb P}
\newcommand{\E}{\mathbb{E}}
\newcommand{\eps}{\varepsilon}
\newcommand{\dd}{\,\mathrm d}
\newcommand{\sref}[2]{\hyperlink{source:#1:#2}{[S#2]}}
\newcommand{\eref}[2]{\hyperlink{extra:#1:#2}{[E#2]}}
% Turn the next switch off for a shorter reading copy. The quotations preserve
% every exactTarget field, including qualifications omitted from a short summary.
\newif\ifshowcatalogue
\showcataloguetrue
\newcommand{\cataloguescope}[1]{\ifshowcatalogue
 \paragraph{Exact scope in the supplied catalogue.}
 \begin{quote}\small #1\end{quote}\fi}

% Standalone wrapper; no external inputs or bibliography files are required.
\fancyhead[L]{\small\sffamily Research notebook: section 207}
\fancyhead[R]{\small 27 September 2026}
\hypersetup{pdftitle={Research attempt 207},pdfauthor={Research notebook}}
\setlength{\parskip}{4pt plus 1pt}
\setlist[itemize]{before=\small\interlinepenalty10000,itemsep=3pt}
\begin{document}
\setcounter{section}{206}
% ================================================================
% problemId: problem.hilbert-smith-conjecture
% source releaseRank: 207; source releaseStatus: open_with_solved_subcases
% exactTarget SHA-256: fc782e193c16deb6757a6cb07adeaffd080cf03f10cabf0cc3dc2727bd1d741b
\clearpage
\section{\texorpdfstring{Hilbert-Smith Conjecture}{Hilbert-Smith Conjecture}}\label{207}
\subsection{Definitions and mathematical statement}
A faithful continuous action of a locally compact Hausdorff group $G$ on a connected topological manifold $M$ is a jointly continuous map $G\times M\to M$ satisfying the action laws, whose only element fixing every point is the identity. The Hilbert--Smith assertion says every such $G$ is a finite-dimensional Lie group. Its standard obstruction formulation excludes faithful actions of
\[
 \Z_p=\varprojlim_m\Z/p^m\Z
\]
by homeomorphisms for every prime $p$. Here $M$ has the usual finite-dimensional manifold conventions, and $\Z_p$ is the compact additive group of $p$-adic integers, not the field $\Q_p$. No differentiability or Lipschitz regularity of the action is assumed.
\par\smallskip\textit{Formulation and concept references:} \sref{207}{1}.
\cataloguescope{Prove that every locally compact topological group acting faithfully on a connected topological manifold is a Lie group; equivalently, prove that Z\_p admits no faithful action by homeomorphisms on a connected manifold.}
\subsection{Short English statement}
Must every locally compact group acting continuously and faithfully on a connected manifold be a Lie group, even when its action is only by homeomorphisms?
% BEGIN RESEARCH ADDITION 207
\subsection{Research attempt}

\paragraph{Current frontier.}
Pardon's formulation source explains both the reduction to $\Z_p$ and
Newman's theorem excluding sufficiently small nontrivial compact Lie
actions on a connected manifold; it also proves the two-dimensional
case.\sref{207}{1}
Pardon's separate three-manifold theorem handles arbitrary continuous
faithful locally compact group actions in dimension three, without
smoothness.\eref{207}{1}
A recent different method is Shelukhin's 2024 Theorem A: no nontrivial
continuous $\Z_p$ action by Hamiltonian homeomorphisms on a \emph{closed
symplectically aspherical} symplectic manifold. Its Floer-theoretic input
belongs to that setting and is not available for arbitrary topological
homeomorphisms.\eref{207}{2}
The attempt below uses no differentiability, invariant symplectic form,
or assumed regularity of an orbit-space metric.

\paragraph{Attack route.}
A regularity route would need control of an invariant metric beyond mere
topological compatibility. An averaging route would need an injectivity
statement not implied by uniform closeness to coordinates. A finite-group
route seems attractive because Newman already gives an order-independent
obstruction, but $\Z_p/p^n\Z_p$ acts naturally on an orbit quotient, not
on the original manifold. The route pursued here makes that last issue
explicit: construct increasingly accurate finite cyclic \emph{almost}
actions on the original manifold and identify the precise repair theorem
that would convert them to genuine actions. The resulting conditional
argument treats noncompact manifolds as well as compact ones.

\paragraph{Attempt.}
\textbf{1. The local obstruction already available from topology.}
We use the following established form of Newman's theorem. If $M$ is
connected, $d$ is a compatible metric, and $U\subset M$ is a nonempty
open set, then there exists $\varepsilon_U>0$ such that a compact Lie
group action with
\[
 d(gx,x)\le\varepsilon_U\quad\text{for all }x\in U\text{ and }g
\]
is trivial on all of $M$. In particular the constant works for every
finite cyclic group, irrespective of its order or whether its action
is faithful. Pardon's source states and proves the version needed
here.\sref{207}{1}
One cannot insert $\Z_p$ into this theorem as though it were a compact
Lie group; that would assume the conclusion.

Fix a prime $p$ and suppose for the purpose of constructing an obstruction
that $a:\Z_p\to\operatorname{Homeo}(M)$ is a faithful continuous action.
Write $a_t(x)$ for the action of $t$. Choose a relatively compact nonempty
open $U$ and a compact set $K_0$ containing $\overline U$. Joint continuity
on compact sets gives
\[
 \sup_{t\in p^m\Z_p}\sup_{x\in K_0}d(a_t x,x)\longrightarrow0
 \quad(m\to\infty). \tag{207.1}
\]
Indeed a neighborhood of $0$ that works uniformly on $K_0$ is obtained
from a finite subcover of $\{0\}\times K_0$, and eventually contains
$p^m\Z_p$. Fix $m$ for which this supremum is less than
$\varepsilon_U/3$. Faithfulness still gives $a_{p^m}\ne\operatorname{id}$.

\textbf{2. Finite cyclic almost actions without changing the phase space.}
For $n\ge1$ and $0\le j<p^n$, define
\[
 \rho_n(j)=a_{p^m j},\qquad
 \rho_n:\Z/p^n\Z\longrightarrow\operatorname{Homeo}(M).
 \tag{207.2}
\]
These are maps of sets with $\rho_n(0)=\operatorname{id}$, not generally
homomorphisms. In a product $j+l$, the integer representative either
stays in range or differs from its residue by $p^n$. Thus the defect is
exactly a carry by $a_{p^{m+n}}$:
\[
 \rho_n(j)\rho_n(l)=
 \begin{cases}
  \rho_n(j+l),&j+l<p^n,\\
  \rho_n(j+l-p^n)a_{p^{m+n}},&j+l\ge p^n.
 \end{cases} \tag{207.3}
\]
For every compact $K\subset M$ it follows that
\[
 \delta_n(K):=
 \max_{j,l\in\Z/p^n\Z}\sup_{x\in K}
 d\bigl(\rho_n(j)\rho_n(l)x,\rho_n(j+l)x\bigr)
 \longrightarrow0. \tag{207.4}
\]
To justify uniformity, consider
$d(a_{t+s}x,a_t x)$ for $t\in p^m\Z_p$, $x\in K$ and
$s=p^{m+n}\to0$. Continuity on this compact parameter set makes the
supremum tend to zero. No quantitative modulus or Lipschitz bound is
needed.

There is additional compactness that an eventual repair theorem should
exploit. All maps in (207.2), and their inverses, lie in the image of the
compact group $p^m\Z_p$. Hence they are equicontinuous on compact sets.
The union of their images of any compact $K$ is contained in the compact
orbit saturation $a(p^m\Z_p\times K)$. This controls both the maps and
where their images can go on a noncompact manifold.

\textbf{3. The precise missing repair principle.}
Here is a sufficient hypothesis, deliberately stated independently of
any assumed $\Z_p$ action.

\emph{Finite cyclic repair hypothesis $\mathrm{FCR}(M,p)$.}
Let $\rho_n:\Z/p^n\Z\to\operatorname{Homeo}(M)$ have identity at zero.
Assume that their maps and inverses, taken together, are equicontinuous
on compact sets; that the union of their images of each compact set has
compact closure; and that the defects (207.4) tend to zero for every
compact set. Then there exist homomorphisms
$\sigma_n:\Z/p^n\Z\to\operatorname{Homeo}(M)$ such that, for every
compact $K$,
\[
 \max_{j\in\Z/p^n\Z}\sup_{x\in K}
 d\bigl(\sigma_n(j)x,\rho_n(j)x\bigr)\longrightarrow0.
 \tag{207.5}
\]
The convergence must be uniform in the \emph{whole} finite cyclic group,
whose order tends to infinity. The hypothesis is not asserted to be a
known theorem or to have been proved in this notebook.

\textbf{Lemma 207.1 (repair implies the obstruction formulation).}
If $\mathrm{FCR}(M,p)$ holds for a connected manifold $M$, there is no
faithful continuous $\Z_p$ action on $M$.

\emph{Proof.}
Apply the repair hypothesis to (207.2), whose hypotheses were verified
in step 2. By (207.1), every $\rho_n(j)$ moves every point of $U$ by
less than $\varepsilon_U/3$. By (207.5) on $K_0$, for all sufficiently
large $n$, every $\sigma_n(j)$ moves those points by less than
$2\varepsilon_U/3$. Each $\sigma_n$ is a genuine finite group action,
so Newman's theorem makes it trivial \emph{on all of $M$}.
For any compact $K$, (207.5) at $j=1$ now gives
\[
 \sup_{x\in K}d(a_{p^m}x,x)
 =\sup_{x\in K}d(\rho_n(1)x,\sigma_n(1)x)
 \longrightarrow0.
\]
The left side is independent of $n$. It is zero on every compact set,
so $a_{p^m}=\operatorname{id}$ globally, contradicting faithfulness.
\hfill$\square$

Consequently, $\mathrm{FCR}(M,p)$ for every connected finite-dimensional
manifold and prime would imply the full target through the standard
locally compact-group reduction.\sref{207}{1}
The proof has not restricted $M$ to be compact: the local Newman
threshold propagates triviality globally, and (207.5) is subsequently
used on arbitrary compact sets. For manifolds with boundary, under
conventions allowing them, a faithful action restricts faithfully to
the interior (a homeomorphism fixing the dense interior fixes the
boundary); the same argument is applied to that interior. A connected
zero-dimensional manifold is a point and its faithful acting group is
trivial.

\textbf{4. Why repairing only one generator would be insufficient.}
The relation that fails in (207.2) is
$\rho_n(1)^{p^n}=a_{p^{m+n}}$, which approaches the identity on compact
sets. Even finding a nearby periodic homeomorphism does not immediately
supply (207.5): one must control all its first $p^n$ powers.
For illustration, on a compact phase space with a metric in which the
original map $f$ is an isometry, a candidate map $g$ with
$\sup_xd(gx,fx)\le\eta$ satisfies only the elementary estimate
\[
 \sup_xd(g^j x,f^j x)\le j\eta.
\]
Thus a bound $\eta_n\to0$ alone is inadequate; this estimate would need
$p^n\eta_n\to0$. Joint continuity gives no such rate for the carry
$a_{p^{m+n}}$. Even a modulus decaying like $1/n$ would be compatible
with continuity but defeat that estimate. This is an obstruction to
this proposed proof method, not a construction of an actual faithful
manifold action with that modulus.

The compact-action invariant metric is easy to construct:
\[
 d_*(x,y)=\sup_{t\in\Z_p}d(a_t x,a_t y)
\]
for a bounded compatible metric $d$. It is a metric, is invariant under
the action, and induces the original topology by compact-parameter
uniform continuity. These facts alone do not make it a Riemannian
metric or give Lipschitz regularity relative to a fixed smooth metric.
An isometry-group theorem with such additional hypotheses therefore
cannot be invoked on $d_*$ without proving them.

\textbf{5. Testing the repair idea in a genuine non-manifold action.}
Take the phase space to be the Cantor group $\Z_p$ itself with its
$p$-adic metric, acting on itself by addition. For $x\in\Z_p$ write
$x=b+p^n y$, $0\le b<p^n$. The translations
$\rho_n(j)x=x+j$, with $j$ represented in $[0,p^n)$, have defect at
most $p^{-n}$. Define an exact action instead by
\[
 \sigma_n(j)(b+p^n y)
 =\bigl((b+j)\bmod p^n\bigr)+p^n y.
 \tag{207.6}
\]
This is a homeomorphism on each clopen residue cylinder and a permutation
of those cylinders. Addition modulo $p^n$ proves the group law. Also
\[
 \sup_{j,x}|\sigma_n(j)x-(x+j)|_p\le p^{-n}.
\]
Hence the desired kind of correction can really occur; it is not merely
a formal request. Yet the original faithful $\Z_p$ action certainly
exists here. The distinction is exactly that Newman's manifold theorem
fails on a Cantor space. Truncated translations by multiples of $p^m$
give nontrivial finite actions all of whose displacements are at most
$p^{-m}$. The conditional proof therefore uses the manifold hypothesis
at an identifiable step and does not incorrectly rule out known
profinite actions on arbitrary compact spaces.

\textbf{6. A second route stops at injectivity, not at averaging.}
On a coordinate region whose relevant small-subgroup orbits remain in
the chart, average its coordinate functions over $p^m\Z_p$ using Haar
probability. Compactness gives continuous averaged coordinates that
are invariant under that subgroup and uniformly close to the original
coordinates as $m\to\infty$. To turn this into an obstruction one
would need strong information about their fibers, and ultimately about
triviality of the action beyond that region.

Uniform closeness does not provide injectivity. Already the real maps
\[
 q_\varepsilon(x)=x+\varepsilon\sin(x/\varepsilon^2)
\]
are within $\varepsilon$ of the identity but have derivative
$1+\varepsilon^{-1}\cos(x/\varepsilon^2)$ of both signs on any fixed
interval when $\varepsilon$ is small. They are not injective there.
This is not claimed to be an average of a hypothetical group action;
it refutes the proposed inference from closeness \emph{alone}.
Degree-one or homotopy-to-identity information also does not imply
injectivity. Additional orbit-fiber topology would still have to be
proved, so this alternative does not fill the repair gap.

\paragraph{Stress test.}
The finite quotients in step 2 were never treated as genuine actions.
The carry formula records their exact failure. The repair hypothesis
includes compact-image control as well as equicontinuity, and uniformity
in all group elements rather than only a generator. Newman's constant
is used only for the repaired finite actions; applying it to the
original $p$-adic group would be circular.

Arzel\`a--Ascoli compactness does not establish (207.5). A convergent
subsequence of some maps can recover a compact group action, which is
precisely the structure one started with, without producing finite-order
homeomorphisms satisfying all the required relations. Neither a
near-identity averaging map nor an invariant compatible metric supplies
that missing algebraic correction. Finally, the scope of the
symplectic theorem has not been enlarged to arbitrary
homeomorphisms.\eref{207}{2}

\Needspace{10\baselineskip}
\paragraph{Outcome.}
\textbf{Outcome: Conditional reduction.}

The attempt gives an explicit family of equicontinuous finite cyclic
almost actions arising from any hypothetical faithful $\Z_p$ action,
computes its defect, and proves that the stated local-uniform repair
principle would settle the obstruction formulation on all connected
manifolds, including noncompact ones. The Cantor model provides an
exact test of the construction and identifies where manifold topology
is indispensable.

No finite cyclic repair theorem with the required variable-order
uniformity has been established. Proving it would itself imply the
full Hilbert--Smith conjecture, so the repair step is a major unresolved
hypothesis, not a consequence of ordinary compactness or continuity.
The invariant-metric and averaging routes were also investigated but
stop at explicit regularity and injectivity gaps. No counterexample to
the original manifold assertion was found.

% Keep the sources heading with a useful initial bibliography block.
\Needspace{8\baselineskip}
% END RESEARCH ADDITION 207
\subsection{Sources}
\begin{itemize}\raggedright
\item[\textnormal{[S1]}]\hypertarget{source:207:1}{} John Pardon, "Totally disconnected groups (not) acting on two-manifolds," manuscript of 6 July 2018 (surveying the Hilbert-Smith Conjecture; published in proceedings on Breadth in Contemporary Topology).
\url{https://www.math.stonybrook.edu/~jpardon/manuscripts/13_hs2d.pdf}.
{\small\textit{Catalogue formulation source.}}
% BEGIN REFERENCE ADDITION 207
\item[\textnormal{[E1]}]\hypertarget{extra:207:1}{}
J. Pardon, \emph{The Hilbert--Smith conjecture for three-manifolds},
Journal of the American Mathematical Society 26 (2013), 879--899,
DOI: 10.1090/S0894-0347-2013-00766-3; main theorem and local reduction.
\url{https://arxiv.org/abs/1112.2324}.
\item[\textnormal{[E2]}]\hypertarget{extra:207:2}{}
E. Shelukhin, \emph{A symplectic Hilbert--Smith conjecture},
arXiv:2403.07195v2 (2024), Theorem A and Section 2, with appendix by
L. Polterovich. The theorem used here assumes a closed symplectically
aspherical manifold and Hamiltonian homeomorphisms.
\url{https://arxiv.org/abs/2403.07195}.
% END REFERENCE ADDITION 207
\end{itemize}

\end{document}
